\chapter{Сравнительное тестирование}\label{ch:results}

\section{Методика тестирования}

Для оценки качества реализованного решения использован стандартный
бенчмарк JSONBench~--- набор аналитических запросов над потоком
JSON-событий из публичного потока сети
Bluesky~\cite{jsonbench}. Бенчмарк содержит около 40~млн событий; в
рамках тестирования использовался срез из 500~000 строк, достаточный
для получения репрезентативных показателей.

Сравниваемые конфигурации подобраны таким образом, чтобы
одновременно показать вклад архитектурного решения, вклад
оптимизаций физических операторов и положение полученной системы
относительно промышленного аналитического движка:

\begin{itemize}
    \item \textbf{Otterbrix (исходная версия).} Сборка ветки
    разработки до начала настоящей работы~--- с документной
    подсистемой хранения, унаследованной от ранних версий проекта.
    Документы хранятся в формате, концептуально близком к BSON в
    MongoDB~\cite{mongodb-bson}: каждый JSON-документ сериализуется
    как неделимое значение и реконструируется построчно при доступе
    к любому из своих полей. Эта конфигурация демонстрирует
    отправную точку~--- производительность движка, основанного на
    документной модели, которую исторически продвигают MongoDB и
    другие документные системы.
    \item \textbf{Otterbrix (новая архитектура).} Сборка с
    реализованным в настоящей работе колоночно-разреженным
    хранением (главы~3--4), но без оптимизаций физических
    операторов из главы~5.
    \item \textbf{Otterbrix (новая архитектура с оптимизациями).}
    Финальная сборка, включающая все оптимизации из главы~5:
    проекцию колонок, типизированные предикаты с пакетным
    копированием, специализированные пути GROUP~BY, исправление
    zone~maps, специализацию служебного идентификатора,
    параллельную запись в side-таблицы в рамках одной
    multi-collection-транзакции.
    \item \textbf{DuckDB~1.1.}~--- встраиваемая колоночная
    аналитическая СУБД~\cite{raasveldt2019}. Выбрана как зрелый
    эталон <<классического>> аналитического движка с многолетней
    историей оптимизации, по архитектурной нише сопоставимый с
    Otterbrix (встраиваемый OLAP-движок в одном процессе с
    приложением). Использован тип \texttt{JSON} с индексами по
    основным полям.
\end{itemize}

Документная система MongoDB рассматривается в работе как
концептуальный прообраз документной модели хранения; численное
сравнение с MongoDB не проводилось напрямую, однако исходная
конфигурация Otterbrix (с собственной документной подсистемой)
отражает тот же класс ограничений, что и BSON-подход MongoDB:
полное чтение документа при доступе к любому полю, отсутствие
кросс-документной колоночной локальности.

Измерения выполнялись на машине с процессором AMD Ryzen~7~5800X
(8~ядер, 16~потоков, 3.8~ГГц), 32~ГБ DDR4, NVMe SSD, под
управлением Ubuntu~24.04 LTS. Все кэши предварительно прогревались
однократным прогоном; фиксировалось среднее время из 5 повторных
прогонов.

\section{Набор запросов}

Использовался стандартный набор из пяти запросов:
\begin{itemize}
    \item \textbf{Q1.} \texttt{SELECT COUNT(*) FROM bluesky}~---
    базовое сканирование.
    \item \textbf{Q2.} \texttt{SELECT kind, COUNT(*) FROM bluesky
    GROUP BY kind}~--- агрегация по закреплённому полю.
    \item \textbf{Q3.} \texttt{SELECT commit.collection, COUNT(*)
    FROM bluesky WHERE commit.operation = 'create' GROUP BY
    commit.collection ORDER BY 2 DESC LIMIT 10}~--- агрегация с
    фильтром.
    \item \textbf{Q4.} \texttt{SELECT commit.record.langs,
    COUNT(*) FROM bluesky GROUP BY commit.record.langs ORDER BY 2
    DESC LIMIT 10}~--- агрегация по разреженному пути.
    \item \textbf{Q5.} \texttt{SELECT did, COUNT(*) FROM bluesky
    WHERE commit.record.reply IS NOT NULL GROUP BY did ORDER BY 2
    DESC LIMIT 10}~--- комбинация фильтра по разреженной колонке и
    агрегации по закреплённой.
\end{itemize}

Q1--Q3 задействуют преимущественно закреплённые поля, что тестирует
путь основной колоночной таблицы. Q4 и Q5 требуют обращения к
разреженным колонкам~--- тест правильности гибридного
сканирования оператором \texttt{scan\_computing\_table} и его
производительности.

\section{Результаты измерений}

Результаты измерений приведены в таблице~\ref{tab:times}. Колонки
расположены в порядке нарастания зрелости решения: исходная
документная конфигурация Otterbrix, новая архитектура без
оптимизаций, новая архитектура с оптимизациями, и эталонный
DuckDB.

\begin{table}[!h]
\caption{Время выполнения запросов JSONBench (с; меньше~--- лучше)}
\label{tab:times}
\centering
\begin{tabular}{|l|r|r|r|r|}\hline
Запрос & Otterbrix (исх.) & Otterbrix (нов.) & Otterbrix (опт.) & DuckDB \\\hline
Q1 & 3.500 & 0.486 & 0.021 & 0.075 \\\hline
Q2 & 4.500 & 1.133 & 0.333 & 0.306 \\\hline
Q3 & 4.600 & 1.498 & 0.273 & 0.262 \\\hline
Q4 & 8.300 & 0.582 & 0.077 & 0.241 \\\hline
Q5 & 8.900 & 0.719 & 0.167 & 0.211 \\\hline
\end{tabular}
\end{table}

Чтобы наглядно показать вклад каждого этапа, в таблице~\ref{tab:speedups}
приведены ускорения относительно исходной конфигурации.

\begin{table}[!h]
\caption{Ускорения относительно исходной конфигурации Otterbrix}
\label{tab:speedups}
\centering
\begin{tabular}{|l|r|r|r|}\hline
Запрос & Архитектура (нов.) & + оптимизации & DuckDB (для сравнения) \\\hline
Q1 & 7.2$\times$  & 167$\times$ & 47$\times$ \\\hline
Q2 & 4.0$\times$  & 13.5$\times$ & 14.7$\times$ \\\hline
Q3 & 3.1$\times$  & 16.9$\times$ & 17.6$\times$ \\\hline
Q4 & 14.3$\times$ & 108$\times$ & 34$\times$ \\\hline
Q5 & 12.4$\times$ & 53$\times$ & 42$\times$ \\\hline
\end{tabular}
\end{table}

\section{Интерпретация результатов}

\subsection{Эффект архитектурного перехода}

Переход от документной модели хранения к колоночно-разреженной
(Otterbrix исх. $\to$ Otterbrix нов.) даёт ускорение от 3$\times$
до 14$\times$ в зависимости от запроса. Наибольший выигрыш
наблюдается на Q4 и Q5~--- запросах с обращением к разреженным
полям: документная подсистема была вынуждена для каждой строки
полностью разбирать BSON-подобную структуру, тогда как новая
архитектура читает только нужные физические колонки или короткие
вспомогательные таблицы. На Q1 (полное сканирование с подсчётом)
ускорение около 7$\times$~--- результат того, что в новой
архитектуре подсчёт строк сводится к обращению к метаданным
row~groups вместо итерации по каждому документу.

Этот результат подтверждает основной тезис работы: документная
модель хранения, лежащая в основе MongoDB и аналогичных систем,
принципиально не приспособлена к аналитической нагрузке, и
переход на колоночное представление даёт качественный, а не
количественный выигрыш.

\subsection{Эффект оптимизаций физических операторов}

Оптимизации главы~5 (Otterbrix нов. $\to$ Otterbrix опт.)
дополнительно ускоряют запросы в 4--23$\times$. Самый яркий эффект
зафиксирован на Q1: ускорение в 23$\times$ относительно неоптимизированной
архитектуры~--- следствие того, что column~pruning превращает
запрос \texttt{COUNT(*)} в обращение к одной служебной колонке без
загрузки остальных. На Q4 ускорение в 7.5$\times$ объясняется
комбинацией column~pruning, типизированного предиката
\texttt{IS~NOT~NULL} и более эффективной реконструкции разреженной
колонки.

В совокупности с архитектурным переходом оптимизации обеспечивают
суммарное ускорение от 13$\times$ (Q2) до 167$\times$ (Q1)
относительно исходной документной конфигурации.

\subsection{Сравнение с DuckDB}

На запросах по закреплённым полям (Q2, Q3) Otterbrix в финальной
конфигурации показывает результаты, практически идентичные DuckDB:
0.333~с против 0.306~с на Q2 (отставание $\sim$9\%) и 0.273~с
против 0.262~с на Q3 (отставание $\sim$4\%). Это означает, что
после применения описанных в работе оптимизаций колоночный путь
исполнения в Otterbrix конкурентоспособен по скорости с одним из
наиболее зрелых открытых аналитических движков.

На запросах по разреженным полям (Q4, Q5) и на полном сканировании
(Q1) Otterbrix \emph{обходит} DuckDB: на Q1 в 3.6$\times$, на Q4
в 3.1$\times$, на Q5 в 1.3$\times$. Причины:

\begin{itemize}
    \item \textbf{Q1.} Финальная реализация \texttt{COUNT(*)}
    обходит загрузку каких-либо колонок данных и считывает счётчик
    из метаданных row~groups, что недоступно в общем пути
    исполнения DuckDB на момент тестирования.
    \item \textbf{Q4 и Q5.} Разреженные поля в Otterbrix хранятся
    в отдельных вспомогательных таблицах, типично содержащих лишь
    долю строк основной таблицы. DuckDB обращается к JSON-полю
    через универсальный экстрактор, что требует полного прохода
    по основной таблице. Преимущество схемы с двухколоночными
    вспомогательными таблицами особенно заметно на нагрузках, где
    редкое поле сочетается с фильтром.
\end{itemize}

На сложных запросах с агрегациями по плотным полям (Q2, Q3) DuckDB
сохраняет небольшое преимущество за счёт зрелой реализации
агрегатных операторов и более агрессивных кодеков сжатия, однако
разрыв укладывается в единицы процентов.

\subsection{Влияние MongoDB-подобной модели на исходную конфигурацию}

Тяжёлые показатели исходной конфигурации (3.5--8.9~с против
0.021--0.333~с в финальной) служат количественной иллюстрацией
известного факта: документная модель хранения, в которой каждое
обращение к полю требует полного разбора документа, проигрывает
колоночному подходу на порядки. Это свойство присуще не конкретной
реализации Otterbrix, а самой документной парадигме~--- те же
ограничения проявляются при использовании MongoDB или PostgreSQL с
JSONB на аналогичных нагрузках. Существенно, что переход на
колоночно-разреженную модель не потребовал отказа от
динамической схемы: пользовательский SQL-интерфейс остаётся прежним,
изменилось только физическое представление данных.

\section{Что достигнуто}

По итогам работы получена практически работоспособная реализация
со следующими характеристиками:
\begin{enumerate}
    \item Производительность на запросах к закреплённым полям
    практически не отличается от DuckDB (отставание единицы
    процентов на Q2 и Q3).
    \item На полном сканировании и на запросах по разреженным
    полям полученная реализация \emph{превосходит} DuckDB в
    1.3--3.6$\times$ за счёт разрешения \texttt{COUNT(*)} через
    метаданные и отдельного хранения редких полей.
    \item Архитектурный переход от документной к
    колоночно-разреженной модели даёт ускорение в 3--14$\times$;
    дополнительные оптимизации физических операторов умножают этот
    выигрыш ещё в 4--23$\times$. Суммарное ускорение по сравнению с
    исходной документной конфигурацией составляет от 13$\times$ до
    167$\times$.
    \item Новый способ хранения JSON в отдельных двухколоночных
    таблицах~--- оригинальный подход, сочетающий преимущества
    колоночного представления и EAV.
    \item Переиспользование существующих механизмов реляционного
    движка: всё хранилище построено на обычных колоночных
    таблицах, что даёт автоматическое распространение всех
    оптимизаций (векторизация, zone~maps, индексы, MVCC, WAL) на
    разреженные колонки без дополнительных усилий.
    \item Не пришлось проектировать отдельную систему хранения для
    разреженных данных; всё опирается на стандартный колоночный
    интерфейс.
\end{enumerate}

\section{Ограничения и направления будущей работы}

Реализация имеет ряд ограничений, которые могут быть сняты в
будущих итерациях:

\begin{itemize}
    \item \textbf{Автоматическое продвижение sparse-поля в pinned.}
    На текущем этапе статус поля фиксируется при
    \texttt{CREATE TABLE} и не пересматривается; ошибочная
    конфигурация (часто используемое поле, оставленное sparse)
    сохраняет повышенную стоимость чтения. Эвристика, которая
    корректирует список pinned-полей на основе реальных запросов и
    статистики заполненности, без нарушения инварианта
    ``main~--- источник истины'' во время миграции, является
    наиболее очевидным направлением развития.
    \item \textbf{Специализированные индексы для side-таблиц.}
    B-дерево по \texttt{\_id} в side-таблицах ускорит резолв
    физических позиций при DELETE и UPDATE на вычисляемые таблицы.
    \item \textbf{Удаление неиспользуемых side-таблиц.} Сжатие и
    выгрузка long-tail side-таблиц с низкой плотностью.
    \item \textbf{Ослабление сериализации запросов на исполнителях.}
    Текущая хеш-маршрутизация одной коллекции на один акторный
    исполнитель является самостоятельным ограничением, не
    связанным с архитектурой вычисляемой схемы; её ослабление
    повысит максимальный параллелизм, и предложенная архитектура к
    этому изменению готова благодаря инвариантам уровня хранилища
    (см.~раздел~\ref{sec:parallel} главы~3).
    \item \textbf{Распределённая репликация.} Шардирование
    side-таблиц по \texttt{\_id} синхронно с main.
    \item \textbf{Догон DuckDB по агрегациям.} Сохранение
    отставания на единицы процентов в Q2 и Q3 указывает на
    потенциал дальнейшей оптимизации хэш-агрегации и кодеков
    сжатия.
\end{itemize}

\chapterconclusion

Результаты сравнительного тестирования подтверждают, что
предложенная архитектура обеспечивает производительность,
конкурентоспособную с признанным промышленным аналитическим
движком DuckDB, и качественно превосходит исходную документную
конфигурацию того же проекта~--- которая сама по себе отражает
модель хранения, исторически продвигаемую MongoDB и аналогичными
документными системами. Ключевое преимущество подхода~--- в
естественной интеграции с существующими механизмами колоночного
реляционного движка: разреженные таблицы являются обычными
таблицами, автоматически получают все оптимизации, правильно
интегрируются с MVCC, WAL и индексами, что и позволяет в ряде
сценариев (Q1, Q4, Q5) обойти DuckDB при на порядок меньшем объёме
кода, специфичного для работы с JSON.
